\cleardoublepage

\section{实验}

\subsection{实验设计}

本节旨在验证 HieraTest 在单元测试生成任务中的有效性，并分析大小模型协同策略对生成性能的影响。实验从生成质量、执行通过情况和代码覆盖能力三个方面展开，主要考察生成测试是否能够通过编译、是否能够通过 Defects4J 触发测试验证，以及是否能够提升目标项目的行覆盖率与分支覆盖率。

\subsubsection{评估实验数据集}

本文计划选取 Defects4J 中具有代表性的开源 Java 项目作为评估数据集。该类数据集包含真实缺陷、历史测试和成熟构建流程，适合用于考察自动生成测试在真实项目中的可用性。后续实验可根据任务规模选择单个项目案例或多个项目集合，以避免实验结果过于依赖个别项目特征。

\subsubsection{评估实验大语言模型}

实验中使用的大模型与小模型分别承担不同职责。大模型用于复杂语义推理、错误修复和高质量候选生成，小模型用于快速生成和局部补充。为了保证实验结论客观，后续应明确记录模型名称、版本、上下文长度、调用方式和提示词配置，并在实验中保持这些变量尽量稳定。

\subsubsection{对比方法}

为了验证 HieraTest 的有效性，实验应与以下几类方法进行对比：仅使用单一大模型的方法、仅使用单一小模型的方法、传统自动化测试生成方法，以及已有的基于大模型的测试生成方法。通过横向对比，可以更清楚地观察协同设计带来的收益。

\subsubsection{评估实验设置}

实验主要采用以下指标：编译通过率用于衡量生成测试是否能够正确构建；运行通过率用于衡量测试是否能够在目标项目中成功执行；行覆盖率和分支覆盖率用于衡量测试对程序行为空间的覆盖程度；此外，还可统计模型调用次数、平均生成耗时和修复迭代轮数，以分析方法的效率与成本。

\subsection{研究问题}

\subsubsection{对比现有单元测试生成方法，HieraTest 的单元测试生成性能}

该问题主要关注 HieraTest 相比单一模型方法是否能够在编译通过率、运行通过率和覆盖率上取得更优结果。若协同机制有效，则应在保持或降低推理成本的前提下，提高生成测试的整体质量。

\subsubsection{HieraTest 中不同模块对单元测试生成性能的影响}

该问题主要分析静态分析、协同生成和协同修复等模块的贡献。通过逐步去除某些模块，可以观察各模块对最终结果的边际作用，从而验证模块设计的必要性。

\subsubsection{HieraTest 在不同大模型上的单元测试生成性能效果}

该问题主要考察 HieraTest 对不同模型组合的适配性。若方法具有良好的通用性，则在不同基础模型上均应表现出相对稳定的提升趋势。

\subsection{实验结果}

\subsubsection{研究问题一实验结果}

本文后续将结合 Defects4J 实验输出，对比分析 HieraTest 与基线方法在编译通过率、运行通过率以及覆盖率上的表现，并讨论其原因。

\subsubsection{研究问题二实验结果}

本文后续将通过消融实验分析各模块对结果的贡献，重点讨论协同生成与协同修复机制是否真正带来稳定增益。

\subsubsection{研究问题三实验结果}

本文后续将比较不同模型组合下的实验结果，分析模型能力差异对生成质量和修复效率的影响。
